\chapter{WolverineDB Formal Specifications}
\label{app:wolverine-db-specifications}

This appendix categorizes the formal specification architecture defined across the 86 specification documents and 139 source files of the \texttt{wolverine-db} cryptographic trust layer.

\section{Specification Taxonomy}
The formal specifications in \texttt{wolverine-db/specs/} are structured into four principal functional domains:

\begin{enumerate}
    \item \textbf{Cryptographic and Merkle State Invariants (Specs 01--22):} Defines deterministic binary serialization, SHA-256 state hashing, and RFC 6962 Merkle tree derivation.
    \item \textbf{Byzantine Fault Tolerance and Consensus Protocols (Specs 23--48):} Specifies a 4-of-5 Byzantine fault tolerance quorum state machine, epoch rotation, and view-change protocols.
    \item \textbf{Change Data Capture (CDC) and Trigger Schemas (Specs 49--68):} Defines append-only database triggers, mutation tracking tables (\texttt{wolverine\_sys.change\_history}), and cryptographic hash chaining.
    \item \textbf{Survivability, Continuous Reconstruction, and Anomaly Detection (Specs 69--86):} Specifies dependency graph reconstruction algorithms and the Sentinel policy gate for automated rollbacks.
\end{enumerate}

\section{Formal Cryptographic Invariants}

\subsection{RFC 6962 Binary Merkle Leaf and Interior Node Derivation}
To prevent second-preimage attacks, leaves and interior nodes utilize distinct prefix domains:
\begin{align}
\label{eq:merkle-leaf}
\text{LeafHash}(D_i) &= \text{SHA-256}(0x00 \mathbin{\Vert} D_i) \\
\label{eq:merkle-interior}
\text{InteriorHash}(L, R) &= \text{SHA-256}(0x01 \mathbin{\Vert} L \mathbin{\Vert} R)
\end{align}
Where $D_i$ is the canonical byte serialization of mutation record $i$, and $\mathbin{\Vert}$ denotes byte concatenation.

\subsection{Append-Only Change History Hash Chain}
Every state modification in \texttt{wolverine\_sys.change\_history} must satisfy the strict cryptographic recurrence relation:
\begin{equation}
\label{eq:hash-chain}
H_k = \text{SHA-256}(H_{k-1} \mathbin{\Vert} \text{Seq}_k \mathbin{\Vert} \text{PayloadDigest}_k)
\end{equation}
Where $H_0 = 0^{256}$ (genesis block), and any retroactive modification to mutation $j < k$ invalidates all subsequent digests $H_m$ ($m \ge j$).

\section{Implementation Status Classification}
\Cref{tab:wdb-spec-status} maps the formal specification domains to their verified implementation status in the codebase.

\begin{table}[htbp]
\centering
\small
\caption{WolverineDB Formal Specifications Implementation Mapping}
\label{tab:wdb-spec-status}
\begin{tabular}{@{}p{3.5cm}p{3.0cm}p{6.8cm}@{}}
\toprule
\textbf{Specification Domain} & \textbf{Status} & \textbf{Codebase Verification Reference} \\ \midrule
\textbf{Merkle Trees \& Signatures} & \textbf{Implemented} & Verified in \texttt{wolverine-db/src/crypto/merkle.ts} and \texttt{approval.ts}. \\ \addlinespace
\textbf{PostgreSQL CDC Triggers} & \textbf{Implemented} & Verified in \texttt{wolverine-db/src/postgres/triggers.sql}. \\ \addlinespace
\textbf{BFT Consensus Engine} & \textbf{Implemented / Tested} & Verified state machine logic in \texttt{src/bft\_hardening/}. \\ \addlinespace
\textbf{Inter-Node Transport} & \textbf{Simulated} & BFT nodes use in-process \texttt{DirectMemoryNetworkTransport}. \\ \addlinespace
\textbf{EVM State Root Anchoring} & \textbf{Simulated} & Merkle root registry uses in-memory \texttt{Map} (\texttt{src/anchors/evm.ts}). \\ \addlinespace
\textbf{Hardware Key Management} & \textbf{Simulated} & Software key derivation without PKCS\#11 HSM bindings. \\ \bottomrule
\end{tabular}
\end{table}
